\documentclass[12pt,a4paper]{report}
\usepackage[T1]{fontenc}
\usepackage[french]{babel}
\usepackage{geometry}
\usepackage{graphicx}
\usepackage{fancyhdr}
\usepackage{listings}
\usepackage{xcolor}
\usepackage{hyperref}
\usepackage{array}
\usepackage{tabularx}
\usepackage{float}
\usepackage{amsmath}
\usepackage{amssymb}

\geometry{left=2.5cm, right=2.5cm, top=2.5cm, bottom=2.5cm}

\pagestyle{fancy}
\fancyhf{}
\rhead{\textit{Projet DevOps - Cloud Infrastructure}}
\lhead{\textit{Rapport Academique}}
\cfoot{\thepage}
\renewcommand{\headrulewidth}{0.4pt}

\lstset{
    basicstyle=\ttfamily\small,
    breaklines=true,
    frame=single,
    tabsize=2,
    showstringspaces=false
}

\title{
    \textbf{RAPPORT DE PROJET DEVOPS} \\[0.5cm]
    \textbf{Infrastructure Cloud CI/CD et Deploiement Kubernetes}
}

\author{
    \textbf{Auteurs:} \\
    LEGNIOUI SIFEDDINE \\
    LOUAT OUSSAMA \\
    AYOUB AFRAOUI \\
    KAWTAR ELBEJJAJI \\[1cm]
    \textit{MEIA - Master Expertises en Informatique Avancee}
}

\date{\today}

\begin{document}

\maketitle

\newpage

\tableofcontents

\newpage

\chapter*{Resume Executif}
\addcontentsline{toc}{chapter}{Resume Executif}

Ce rapport presente une solution complete d'infrastructure cloud et DevOps, combinant les meilleures pratiques modernes en matiere de conteneurisation, orchestration et infrastructure-as-code.

Le projet implemente une pipeline CI/CD entierement automatisee, permettant le deploiement rapide et fiable d'applications web, tout en assurant la scalabilite, la haute disponibilite et la maintenabilite de l'infrastructure.

\textbf{Principaux objectifs realises:}
\begin{enumerate}
    \item Developpement d'une application web Python moderne
    \item Conteneurisation avec Docker pour portabilite
    \item Infrastructure-as-Code avec Terraform
    \item Automatisation avec Ansible pour configuration/deploiement
    \item Orchestration Kubernetes pour haute disponibilite
    \item Pipeline CI/CD integre pour deploiement continu
    \item Documentation complete et reproductibilite
\end{enumerate}

\textbf{Technologies cles:}
\begin{itemize}
    \item \textbf{Application:} Python Flask
    \item \textbf{Containerisation:} Docker
    \item \textbf{Infrastructure:} Terraform (IaC)
    \item \textbf{Configuration:} Ansible
    \item \textbf{Orchestration:} Kubernetes
    \item \textbf{CI/CD:} GitHub Actions
\end{itemize}

\newpage

\chapter{Introduction}

\section{Contexte du Projet}

L'infrastructure informatique moderne exige une approche holistique combinant plusieurs disciplines : development, operations, et security. Le paradigme DevOps emerge comme solution pour eliminer les silos organisationnels et accelerer le time-to-market tout en assurant la qualite et la securite.

Ce projet s'inscrit dans cette demarche en implementant une approche complete d'automatisation et d'orchestration sur cloud infrastructure.

\section{Objectifs du Projet}

Les objectifs principaux de ce projet sont:

\begin{itemize}
    \item Concevoir et implementer une application web scalable et maintenable
    \item Mettre en place une infrastructure cloud automatisee et reproductible
    \item Implementer une pipeline CI/CD pour automatiser les tests et deploiements
    \item Assurer la haute disponibilite et la resilience de l'application
    \item Documenter toutes les etapes pour faciliter la reproduction et la maintenance
    \item Demontrer les meilleures pratiques DevOps et Cloud Native
\end{itemize}

\section{Perimetre de la Solution}

Cette solution couvre:

\begin{enumerate}
    \item Application Flask Python simple mais fonctionnelle
    \item Dockerfile pour packaging de l'application
    \item Scripts Terraform pour provision cloud
    \item Playbooks Ansible pour configuration serveurs
    \item Manifests Kubernetes YAML pour deploiement
    \item Configuration CI/CD pour automatisation
    \item Guide complet de reproduction et maintenance
\end{enumerate}

\newpage

\chapter{Architecture et Design}

\section{Vue d'ensemble}

L'architecture proposee suit une approche microservices avec orchestration Kubernetes, permettant scalabilite horizontale et haute disponibilite.

\subsection{Composants Principaux}

\begin{itemize}
    \item Application dockerisee et scalable
    \item Kubernetes pour gestion des conteneurs
    \item Cloud provider gerant par Terraform
    \item Ansible pour configuration initiale
    \item Pipeline CI/CD automatisee
\end{itemize}

\section{Flux de Deploiement}

\begin{enumerate}
    \item Le developpeur commit du code sur le repository Git
    \item La pipeline CI/CD est declenchee automatiquement
    \item Tests unitaires et d'integration sont executes
    \item L'image Docker est construite et poussee vers registry
    \item Kubernetes applique les nouveaux manifests
    \item L'application est mise a jour en production
\end{enumerate}

\newpage

\chapter{Application Web - Flask}

\section{Vue d'ensemble}

L'application est une web application Flask Python simple, servant comme cas d'usage demontrant les capacites de l'infrastructure DevOps.

\section{Code Source}

\subsection{Fichier Principal - app.py}

\begin{lstlisting}
from flask import Flask, render_template
import os

app = Flask(__name__)

@app.route('/')
def home():
    return render_template('index.html')

@app.route('/health')
def health():
    return {'status': 'healthy'}, 200

@app.route('/version')
def version():
    return {'version': '1.0.0'}, 200

if __name__ == '__main__':
    app.run(host='0.0.0.0', port=5000)
\end{lstlisting}

\subsection{Dependances - requirements.txt}

\begin{lstlisting}
Flask==2.3.0
Werkzeug==2.3.0
\end{lstlisting}

\section{Specifications Techniques}

\begin{table}[H]
\centering
\begin{tabularx}{\linewidth}{|l|X|}
\hline
\textbf{Parametre} & \textbf{Valeur} \\
\hline
Langage & Python 3.9+ \\
Framework & Flask 2.3.0 \\
Port Application & 5000 \\
Endpoints & /, /health, /version \\
\hline
\end{tabularx}
\end{table}

\newpage

\chapter{Containerisation - Docker}

\section{Concept et Avantages}

Docker est une plateforme de containerisation qui encapsule l'application et toutes ses dependances dans une image portable.

\subsection{Avantages de Docker}

\begin{itemize}
    \item Portabilite: ``Build once, run anywhere''
    \item Isolation des applications
    \item Efficacite comparee aux VMs
    \item Rapidite de demarrage
    \item Scalabilite facile
\end{itemize}

\section{Dockerfile}

\begin{lstlisting}
FROM python:3.9-slim
WORKDIR /app
COPY requirements.txt .
RUN pip install --no-cache-dir -r requirements.txt
COPY app.py .
COPY templates/ templates/
COPY static/ static/
EXPOSE 5000
CMD ["python", "app.py"]
\end{lstlisting}

\section{Commandes Docker Essentielles}

\begin{table}[H]
\centering
\begin{tabularx}{\linewidth}{|l|X|}
\hline
\textbf{Commande} & \textbf{Description} \\
\hline
docker build -t app:1.0 . & Construire image \\
docker run -p 5000:5000 app & Executer conteneur \\
docker push registry/app:1.0 & Pusher vers registry \\
docker ps & Lister conteneurs actifs \\
docker logs <id> & Voir logs \\
\hline
\end{tabularx}
\end{table}

\newpage

\chapter{Infrastructure as Code - Terraform}

\section{Introduction a Terraform}

Terraform est un outil Infrastructure as Code permettant de:
\begin{itemize}
    \item Definir infrastructure en code declaratif
    \item Versionner et controler infrastructure
    \item Deployer infrastructure sur multiples environnements
    \item Faciliter destruction et recreation
\end{itemize}

\section{Structure Terraform}

Les fichiers principaux incluent:

\begin{itemize}
    \item providers.tf - Configuration des providers
    \item variables.tf - Variables d'entree
    \item main.tf - Ressources principales
    \item outputs.tf - Sorties d'execution
\end{itemize}

\section{Flux Terraform}

\begin{enumerate}
    \item terraform init - Initialiser
    \item terraform plan - Afficher changements
    \item terraform apply - Appliquer changements
    \item terraform destroy - Detruire ressources
\end{enumerate}

\newpage

\chapter{Configuration - Ansible}

\section{Vue d'ensemble}

Ansible est un outil de configuration management permettant:
\begin{itemize}
    \item Automatiser deploiement et configuration
    \item Gerer infrastructures complexes
    \item Rester agentless (SSH uniquement)
    \item Ecrire playbooks lisibles
\end{itemize}

\section{Inventaire}

\begin{lstlisting}
[all]
app_server_1 ansible_host=10.0.1.100
app_server_2 ansible_host=10.0.2.100

[app]
app_server_1
app_server_2
\end{lstlisting}

\section{Execution des Playbooks}

\begin{lstlisting}
ansible-playbook playbook.yml
ansible-playbook playbook.yml -i inventory.ini -l app
ansible-playbook playbook.yml --check
\end{lstlisting}

\newpage

\chapter{Orchestration - Kubernetes}

\section{Vue d'ensemble}

Kubernetes est un systeme d'orchestration permettant:
\begin{itemize}
    \item Automatiser deploiement des conteneurs
    \item Scalabilite horizontale
    \item Haute disponibilite et resilience
    \item Load balancing
    \item Gestion des ressources
\end{itemize}

\section{Concepts Cles}

\subsection{Pod}
Unite minimale de deploiement contenant un ou plusieurs conteneurs.

\subsection{Deployment}
Gere replication et updates de Pods.

\subsection{Service}
Expose Pods de maniere stable via IP et DNS.

\subsection{Namespace}
Fournit isolation logique entre ressources.

\section{Deployment Manifest}

\begin{lstlisting}
apiVersion: apps/v1
kind: Deployment
metadata:
  name: devops-app
  namespace: production
spec:
  replicas: 3
  selector:
    matchLabels:
      app: devops-app
  template:
    metadata:
      labels:
        app: devops-app
    spec:
      containers:
      - name: app
        image: devops-app:latest
        ports:
        - containerPort: 5000
        resources:
          requests:
            cpu: 100m
            memory: 128Mi
\end{lstlisting}

\section{Service Manifest}

\begin{lstlisting}
apiVersion: v1
kind: Service
metadata:
  name: devops-app
  namespace: production
spec:
  type: LoadBalancer
  selector:
    app: devops-app
  ports:
  - protocol: TCP
    port: 80
    targetPort: 5000
\end{lstlisting}

\section{Commandes kubectl}

\begin{table}[H]
\centering
\begin{tabularx}{\linewidth}{|l|X|}
\hline
\textbf{Commande} & \textbf{Description} \\
\hline
kubectl apply -f deployment.yaml & Creer ressources \\
kubectl get pods & Lister pods \\
kubectl logs <pod> & Voir logs \\
kubectl describe pod <pod> & Details \\
kubectl scale deployment app --replicas=5 & Scaler \\
\hline
\end{tabularx}
\end{table}

\newpage

\chapter{Pipeline CI/CD}

\section{Concept et Architecture}

Une pipeline CI/CD automatise:
\begin{enumerate}
    \item Continuous Integration - Tests automatiques
    \item Continuous Delivery - Releases pretes
    \item Continuous Deployment - Deploiement automatique
\end{enumerate}

\section{Flux Typique}

\begin{enumerate}
    \item Code pushed au repository
    \item Tests automatiques executes
    \item Image Docker construite
    \item Image poussee vers registry
    \item Application deployee
    \item Tests de sante executes
    \item Equipes notifiees
\end{enumerate}

\section{GitHub Actions}

\begin{lstlisting}
name: CI/CD Pipeline

on:
  push:
    branches: [ main ]

jobs:
  test:
    runs-on: ubuntu-latest
    steps:
      - uses: actions/checkout@v3
      - name: Run tests
        run: pytest tests/

  build:
    needs: test
    runs-on: ubuntu-latest
    steps:
      - uses: actions/checkout@v3
      - name: Build image
        run: docker build -t app:latest .
\end{lstlisting}

\newpage

\chapter{Securite et Gestion des Secrets}

\section{Principes de Securite}

Les secrets ne doivent JAMAIS etre:
\begin{itemize}
    \item Directement dans le code
    \item Commits dans version control
    \item Loggees
    \item Visibles en plain text
\end{itemize}

\section{Gestion de Secrets}

\subsection{Variables d'Environnement}

\begin{lstlisting}
export DB_PASSWORD="secret"
export API_KEY="secret"
python app.py
\end{lstlisting}

\subsection{Kubernetes Secrets}

\begin{lstlisting}
apiVersion: v1
kind: Secret
metadata:
  name: app-secrets
type: Opaque
stringData:
  db_password: "secure_password"
  api_key: "secure_key"
\end{lstlisting}

\section{Bonnes Pratiques}

\begin{enumerate}
    \item Least Privilege - Permissions minimales
    \item Defense in Depth - Multiples couches
    \item Regular Updates - Maintien a jour
    \item Monitoring - Audit trails detailles
    \item Network Isolation - Segmentation
    \item RBAC - Role-Based Access Control
\end{enumerate}

\newpage

\chapter{Monitoring et Logging}

\section{Types de Monitoring}

\begin{itemize}
    \item Infrastructure - CPU, Memoire, Disque, Reseau
    \item Application - Temps reponse, Erreurs
    \item Business - Utilisateurs actifs
    \item Security - Tentatives non-autorisees
\end{itemize}

\section{Structured Logging}

Les logs doivent etre structures en JSON pour faciliter parsing et analyse.

\section{Kubernetes Logging}

\begin{lstlisting}
kubectl logs deployment/devops-app
kubectl logs -f deployment/devops-app
kubectl logs deployment/devops-app --tail=100
\end{lstlisting}

\section{Stack ELK}

Elasticsearch, Logstash, Kibana pour logging centralise.

\newpage

\chapter{Documentation et Deploiement}

\section{Prerequis}

\begin{itemize}
    \item Docker et Docker Compose (v20.10+)
    \item Terraform (v1.0+)
    \item Ansible (v2.10+)
    \item kubectl (v1.24+)
    \item Git
\end{itemize}

\section{Installation Locale}

\begin{lstlisting}
git clone https://github.com/org/devops-project.git
cd devops-project
docker-compose up -d
curl http://localhost:5000/health
\end{lstlisting}

\section{Deploiement Cloud}

\begin{lstlisting}
cd terraform
terraform init
terraform plan
terraform apply

aws eks update-kubeconfig --region us-east-1

kubectl apply -f kubernetes/
kubectl rollout status deployment/devops-app
\end{lstlisting}

\newpage

\chapter{Troubleshooting et Maintenance}

\section{Problemes Docker}

\subsection{Conteneur ne demarre pas}

\begin{lstlisting}
docker logs <container_id>
docker inspect <container_id>
docker run -it <image> /bin/bash
\end{lstlisting}

\section{Problemes Kubernetes}

\subsection{Pod ne demarre pas}

\begin{lstlisting}
kubectl describe pod <pod_name>
kubectl logs <pod_name>
kubectl get events
\end{lstlisting}

\section{Maintenance Reguliere}

\begin{itemize}
    \item Updates des base images mensuels
    \item Verification vulnerabilites
    \item Mise a jour dependances
    \item Audit acces secrets
    \item Tests de disaster recovery
\end{itemize}

\section{Operational Runbooks}

\subsection{Scaler Application}

\begin{lstlisting}
kubectl scale deployment devops-app --replicas=5
kubectl get deployment devops-app
\end{lstlisting}

\subsection{Rollback}

\begin{lstlisting}
kubectl rollout history deployment/devops-app
kubectl rollout undo deployment/devops-app
\end{lstlisting}

\newpage

\chapter{Conclusions et Recommandations}

\section{Resume des Realisations}

Ce projet a demontre la mise en place reussie d'une infrastructure complete DevOps incluant:

\begin{enumerate}
    \item Application web containerisee et scalable
    \item Infrastructure cloud reproductible via Terraform
    \item Automatisation complete des deploiements
    \item Orchestration Kubernetes pour haute disponibilite
    \item Pipeline CI/CD entierement automatisee
    \item Documentation exhaustive
\end{enumerate}

\section{Points Cles de la Solution}

\subsection{Architecture Moderne}

\begin{itemize}
    \item Microservices et conteneurs
    \item Stateless applications
    \item Infrastructure as Code
    \item Immutable infrastructure
\end{itemize}

\subsection{Automatisation Maximale}

\begin{itemize}
    \item Provisioning infrastructure
    \item Configuration servers
    \item Build et tests
    \item Deploiement application
    \item Scalabilite
    \item Monitoring
\end{itemize}

\section{Recommandations Futures}

\subsection{Court Terme (1-3 mois)}

\begin{enumerate}
    \item Monitoring avec Prometheus et Grafana
    \item Logging centralise avec ELK
    \item Alerting via PagerDuty
    \item Security scanning en CI/CD
    \item Tests de charge
\end{enumerate}

\subsection{Moyen Terme (3-6 mois)}

\begin{enumerate}
    \item Service mesh (Istio)
    \item Policy management (OPA)
    \item Disaster recovery testing
    \item A/B testing framework
    \item Cost optimization
\end{enumerate}

\subsection{Long Terme (6+ mois)}

\begin{enumerate}
    \item Multi-region expansion
    \item Machine Learning ops
    \item Cost allocation
    \item Multi-cloud strategy
    \item Serverless components
\end{enumerate}

\section{Lessons Learned}

\begin{itemize}
    \item Infrastructure as Code cruciale pour reproductibilite
    \item Automatisation minimise erreurs humaines
    \item Testing multiples layers obligatoire
    \item Monitoring depuis jour 1
    \item Documentation vitale
    \item Securite baseline importante
    \item Adoption progressive mieux que big-bang
\end{itemize}

\section{Conclusion Finale}

Ce projet demontre comment les technologies modernes peuvent etre combinees pour creer une infrastructure robuste, scalable et maintenable.

Les pratiques DevOps et Cloud-Native implementees etablissent une foundation solide pour:

\begin{itemize}
    \item Accelerer time-to-market
    \item Ameliorer reliability
    \item Reduire operational overhead
    \item Faciliter collaboration Dev-Ops
    \item Supporter croissance future
\end{itemize}

La cle du succes reside dans l'automatisation complete, la documentation claire, et l'adoption continue de meilleures pratiques.

\newpage

\appendix

\chapter{Ressources Additionnelles}

Ressources recommandees:

\begin{itemize}
    \item Kubernetes: https://kubernetes.io/docs/
    \item Terraform: https://registry.terraform.io/
    \item Ansible: https://docs.ansible.com/
    \item Docker: https://docs.docker.com/
    \item GitHub Actions: https://docs.github.com/en/actions
\end{itemize}

\chapter{Auteurs du Projet}

\begin{itemize}
    \item LEGNIOUI SIFEDDINE
    \item LOUAT OUSSAMA
    \item AYOUB AFRAOUI
    \item KAWTAR ELBEJJAJI
\end{itemize}

Module: MEIA - Master Expertises en Informatique Avancee

\end{document}
